\documentclass[12pt]{article}
\usepackage[T1]{fontenc}
\usepackage[utf8]{inputenc}
\usepackage{lmodern}
\usepackage{textcomp}
\usepackage{enumitem}
\usepackage{geometry}
\geometry{margin=1in}
\begin{document}
\begin{center}
Answer Key - Political Science - Undergraduate Level Assessment 25 \
\small (Answers are ordered exactly as in the question paper.)
\end{center}

\section*{Multiple Choice}
\begin{enumerate}[label=\arabic*.]
\item \textbf{Answer:} D
\\ \textit{Options:} A) A defendant who wants to know what witnesses the plaintiff plans to call.; B) Former presidents who wish to send instructions to the justices who they appointed.; C) A senator who wants to impeach a Supreme Court justice.; D) Companies that are not involved in a court case but wish to affect the outcome.
\\ \textit{Note:} Companies that are not involved in a court case may file an amicus curiae brief to present additional information, perspectives, or arguments that could influence the court's decision, reflecting their vested interests in the case's outcome.
\item \textbf{Answer:} C
\\ \textit{Options:} A) One man, one vote; B) Separate but equal; C) Judicial review; D) Right to privacy
\\ \textit{Note:} The Supreme Court's decision in Marbury v. Madison established the principle of judicial review by affirming the Court's authority to invalidate laws and executive actions that are deemed unconstitutional, thereby solidifying the judiciary's role as a check on legislative and executive powers.
\item \textbf{Answer:} D
\\ \textit{Options:} A) The executive and legislative branches working on legislation together; B) The federal government granting power over a policy area to the states; C) Governments working with businesses to address an issue; D) State and federal governments working on the same issue
\\ \textit{Note:} Cooperative federalism is characterized by the collaboration and shared responsibilities between state and federal governments in addressing common issues, reflecting a dynamic partnership rather than a strictly divided authority.
\item \textbf{Answer:} B
\\ \textit{Options:} A) vice president; B) president's chief of staff; C) national security advisor; D) chair of the Federal Reserve Board
\\ \textit{Note:} The president's chief of staff plays a crucial role in managing the flow of information and access between the cabinet secretaries and the president, effectively controlling how and when they can interact.
\item \textbf{Answer:} C
\\ \textit{Options:} A) voters select the winner by caucus instead of by individual ballots; B) the election results are not binding; C) any registered voter may participate, regardless of party affiliation; D) voters may register to vote on the day of the election
\\ \textit{Note:} In an open primary election, the key characteristic is that it allows any registered voter to participate in the primary election process without being restricted by their party affiliation, unlike closed primaries which require voters to be registered with a specific party.
\end{enumerate}

\section*{True / False}
\begin{enumerate}[label=\arabic*.]
\setcounter{enumi}{5}
\item \textbf{Answer:} True
\\ \textit{Options:} A) True; B) False
\\ \textit{Note:} The Supreme Court decision in Roe v. Wade primarily focused on the right to privacy regarding a woman's right to choose an abortion, which is rooted in the Fourteenth Amendment, and did not address any issues pertaining to the First Amendment.
\item \textbf{Answer:} False
\\ \textit{Options:} A) True; B) False
\\ \textit{Note:} The Liberal approach to US foreign policy prioritizes alliances and diplomacy but emphasizes free trade and international cooperation rather than protectionist measures.
\item \textbf{Answer:} True
\\ \textit{Options:} A) True; B) False
\\ \textit{Note:} The statement is true because 'iron triangles' refer to the stable relationships among congressional committees, bureaucratic agencies, and interest groups, while courts operate independently and do not engage in these direct policy-making alliances.
\item \textbf{Answer:} False
\\ \textit{Options:} A) True; B) False
\\ \textit{Note:} Voter behavior is influenced by a multitude of factors including socioeconomic status, education, race, political ideology, and specific issues at stake in the election, making gender alone insufficient to determine voting outcomes.
\item \textbf{Answer:} False
\\ \textit{Options:} A) True; B) False
\\ \textit{Note:} Realists argue that American exceptionalism fosters a unilateral approach to foreign policy that undermines international cooperation and the feasibility of effective world government, as it prioritizes national interests over collective global governance.
\end{enumerate}

\section*{Long Form Response}
\begin{enumerate}[label=\arabic*.]
\setcounter{enumi}{10}
\item \textbf{Answer:} Throughout the twentieth century, the concept of health as a security issue has significantly evolved, reflecting a transition from a predominant focus on military threats to the recognition of pandemic diseases as critical security concerns. Initially, security analysts were primarily preoccupied with the dangers posed by nuclear confrontation and conventional warfare, which shaped the geopolitical landscape of the Cold War era. However, as the threat of military conflict diminished, attention shifted toward public health crises, particularly the emergence of infectious diseases such as HIV/AIDS, SARS, and more recently, COVID-19. This shift highlights a growing awareness of how global interconnectedness can facilitate the rapid spread of diseases, underscoring the necessity for comprehensive health security measures. Consequently, the understanding of security has expanded to encompass health-related issues, prompting nations to prioritize pandemic preparedness alongside traditional defense strategies.
\\ \textit{Note:} correct as it identifies the significant shift in the perception of health as a security issue from military threats during the Cold War to the growing recognition of pandemic diseases as major security concerns, reflecting the evolving nature of global threats over the twentieth century.
\item \textbf{Answer:} In the technical discourse on cyber-security, key concepts include threat modeling, vulnerability assessment, risk management, and incident response, all aimed at protecting information systems from cyber threats. However, business networks often fall outside this framework because they are typically viewed as operational entities focused on efficiency and profitability rather than as security-centric systems. This perspective can lead to a disconnect between cyber-security protocols and the practices of business networks, which may prioritize immediate business objectives over long-term security investments. Additionally, the complexity and diversity of business networks make it challenging to apply uniform cyber-security measures, resulting in a tendency to treat them as separate from the broader cyber-security landscape. Consequently, this separation can leave business networks vulnerable to cyber threats, highlighting the need for an integrated approach that considers both security and operational efficiency.
\\ \textit{Note:} correct because it accurately identifies essential cyber-security concepts while highlighting the operational focus of business networks, which can result in a lack of alignment with security measures.
\end{enumerate}

\vfill
\noindent\textit{End of Answer Key}
\end{document}